\section{Latar Belakang}

Setelah pandemi COVID-19 berakhir, tren olahraga lari di Indonesia mengalami peningkatan yang signifikan. Olahraga ini menjadi salah satu aktivitas favorit masyarakat karena dinilai dapat menjaga kesehatan fisik sekaligus memberikan kesempatan untuk bersosialisasi di ruang terbuka. Menurut data dari KalenderLari, terdapat 257 event lari yang diadakan di Indonesia sepanjang tahun 2024, meningkat lebih dari 60\% dibandingkan tahun sebelumnya. Peningkatan ini juga terlihat dari laporan Viva.co.id, yang menyatakan bahwa kesadaran masyarakat terhadap gaya hidup sehat telah membuat olahraga lari semakin diminati \cite{sumiyati}. Selain itu, pelatih balap lari Andriyanto menyebutkan bahwa jumlah pelari meningkat tajam pasca pandemi, mencerminkan antusiasme masyarakat terhadap olahraga ini \cite{tazkiarh}.

Peningkatan tren lari juga memberikan dampak signifikan pada industri terkait, seperti penjualan perlengkapan olahraga. Laporan dari Economy.Okezone.com menunjukkan bahwa setelah pandemi COVID-19, penjualan sepatu lari meningkat hingga 200\% karena banyak orang yang semakin sadar akan pentingnya kesehatan. Hal serupa juga dikonfirmasi oleh Okezone.com, yang mencatat lonjakan permintaan terhadap pakaian dan aksesori olahraga lari sebagai akibat dari tren ini \cite{muhammadakbar}.

Di sisi lain, pembelian tiket untuk acara olahraga, termasuk lari, kini lebih banyak dilakukan secara online. Platform manajemen event seperti Loket.com dan TixFly mencatat peningkatan signifikan dalam penjualan tiket digital, yang memperlihatkan pergeseran preferensi konsumen ke sistem yang lebih praktis dan efisien. Dalam artikel yang diterbitkan oleh Media Indonesia, platform digital kini menjadi pilihan utama penyelenggara event karena dapat memangkas biaya operasional serta memberikan kenyamanan lebih bagi pembeli tiket \cite{despian}.

Pengembangan aplikasi e-ticketing untuk event lari menjadi hal yang penting. Aplikasi ini tidak hanya mempermudah proses pembelian tiket, tetapi juga mendukung pengelolaan acara yang lebih efisien. Dengan fitur seperti pembuatan event, pengelolaan tiket, dan laporan penjualan yang terintegrasi, e-ticketing memungkinkan penyelenggara event untuk bekerja lebih efektif. Bagi pembeli tiket, aplikasi ini memberikan kemudahan dalam mengakses informasi acara, memilih jenis tiket, hingga melakukan pembayaran secara online kapan saja dan di mana saja. Seperti yang dijelaskan oleh Artika Surniandari dan Haryani \cite{surniandari2017}, kualitas informasi, kualitas sistem, dan kualitas pelayanan sangat berpengaruh terhadap kepuasan dan loyalitas pengguna dalam sistem e-ticketing.

Selain itu, masalah autentikasi menjadi salah satu perhatian utama dalam pengembangan aplikasi e-ticketing. Studi dari NordPass \cite{nordpass2023} menunjukkan bahwa sekitar 73\% pengguna sering lupa kata sandi mereka, yang menyebabkan frustrasi dan menghambat pengalaman pengguna. Dalam survei yang dilakukan oleh Statista pada 2023 \cite{statista2022}, lebih dari 71\% pengguna lebih memilih menggunakan login Google atau OAuth karena lebih cepat dan aman dibandingkan mengingat kata sandi untuk setiap akun yang mereka miliki. Oleh karena itu, aplikasi e-ticketing ini dirancang dengan sistem passwordless authentication menggunakan OAuth Google. Fitur ini tidak hanya meningkatkan kenyamanan pengguna, tetapi juga memperkuat keamanan dengan mengurangi risiko serangan berbasis kata sandi, seperti phishing atau credential stuffing.

Oleh karena itu, pengembangan aplikasi e-ticketing untuk event lari menjadi solusi yang tidak hanya menjawab kebutuhan pasar yang terus berkembang, tetapi juga memberikan pengalaman yang praktis, aman, dan modern. Dengan fitur seperti passwordless authentication menggunakan OAuth Google dan integrasi webhook untuk multi-channel pembayaran, aplikasi ini diharapkan dapat memberikan nilai tambah yang signifikan, baik bagi penyelenggara event maupun peserta lari.